\section{相关工作}

\subsection{秘密警察与规训的现代研究}福柯（Foucault）\cite{foucault1977}把惩罚方式的转变解释为
规训权力的技术化，全景敞视（panopticon）恰恰是「被看见的权力」与「看不见的监督者」的合体：
监视者不需要在场，被监视者必须假设其在场。这一视角给本文的启发是：**监视作为一种治理术，
其效力来自信息的不对称，而不是暴力本身**。本文的Fig.3 那组数字正是对这一论断的一次反向检验——
在本书的图谱里，行使监视的机构（厂卫系 29 人）在结构上并不比普通人物更中心，
它的权力是叙事赋权的，不是关系位置赋予的。需要指出的是，福柯研究的是近代欧洲的狱舍与学校，
明代的检校、番役与诏狱与之在制度载体上不可通约，本文只借用其「可见性—不可见性」的辨析框架，
不做制度层面的因果类比。

\subsection{偏好伪造与集体说谎}库兰（Kuran）\cite{kuran1995}提出偏好伪造（preference falsification）：
个体私下持有与公开表达不同的判断，并且因为害怕孤立即沉默，最终使群体走向一个谁都不信的共识。
这一机制与本文所说的「告密—黑名单」结构最直接相关：明代政治里最有效的攻击手段往往不是证据，
而是让对手先自证、再把自证变成罪状。本文数据中「政治攻讦」端点里有 86.0\% 属跨阵营配对（见结果节的跨阵营统计），
说明攻讦主要发生在对立阵营之间而非阵营内部，这一点与偏好伪造预测的「公开表达趋同、私下判断分歧」并不矛盾，
但偏好伪造是微观动机理论，本文是语料结构测量，二者只能在方向上呼应，不能被本文的数字验证。

\subsection{公共领域与古代帝国的的信息公开}哈贝马斯（Habermas）\cite{habermas1989}的公共 sphere 讨论的是
现代印刷公众的形成条件；把它套到明代政治上要特别小心。本书的语料里，公开渠道词的逐字命中数是
奏疏 202、弹劾 179、上疏 88，而隐秘渠道是锦衣卫 207、东厂 122、特务 86、妖书 31、告密 4、黑名单 5
（表 II）。两组量级相当，但**性质完全不同**：奏疏与弹劾是制度化的、署名可追的、可存档的；
告密与黑名单是不署名的、不可存档的、只服务于当场的。把这两组一起放进词频表，恰恰说明本书写的是
一个公开程序与隐秘机制并存的双通道系统，而不是单通道的舆论场。本文因此不使用「公共领域」一词，
改用「公开／隐秘双通道」这一操作化说法。

\subsection{抵抗的隐秘文本}斯科特（Scott）\cite{scott1990}区分公开的台词与隐藏的 transcript，
认为弱者的日常抵抗总在正式记录之外。本文的抽取数据给出一个反向的意外：本书里被正式记录、被反复书写的信息
恰恰是公开程序（弹劾、奏疏），而真正构成杀伤的隐秘动作（告密、黑名单、谣言）词频只有个位数到几十次。
这不能用来否定斯科特——它是本书的选择，不是明代的实况——但它提示：**通俗史著的「史料化」过程是选择性
的**，我们把一个王朝的隐秘 transcripts 读成正式文档的密度，本身就是这本书编辑判断的产物。

\subsection{数字人文中的语料计算}贝里（Berry）\cite{berry2011}把大规模数字化语料的统计转向列为数字人文的
第一条主线，同时指出这一转向的固有争议：大规模统计会系统性地偏爱高频词汇与常见结构。
本文的Fig.2 正是这一争议在
单一语料上的一次内向检验。与既有数字人文工作相比，本文的差别在于**不把词频当作内容**，而是把词频
当作需要被对照的因变量：先看词频（谁被反复写），再看结构位置（谁真正成了网络节点），两者不一致的地方
恰恰是本文的主结论。

\subsection{明代政治史的既有结论}莫特与崔瑞德主编的《剑桥中国明代史》\cite{mote1988}提供了制度—结构路径的
标准参照，黄仁宇《万历十五年》\cite{huang1981}则提供了把孤立年份读成结构征兆的微观做法。
本文与二者的关系不是替代，而是补充：它们读的是档案里的明代，本文读的是一个现代人写给大众的明代，
两者的差别本身就是研究对象的一部分。
